% META: {"id": "TSPE-LOG-EX-003", "chapitre": "TSPE-LOGARITHME", "type_objet": "exercice", "capacites": ["TSPE-LOG-C2"], "capacites_codes": ["C2"], "methodes": ["M1"], "parcours": 2, "competences": ["modeliser", "calculer"], "duree_min": 15, "mode_creation": "ex_nihilo", "sources_inspiration": [], "fichier_tex": "chapitres/TSPE-LOGARITHME/exercices/TSPE-LOG-EX-003.tex", "corrige_tex": "chapitres/TSPE-LOGARITHME/corriges/TSPE-LOG-CO-003.tex", "status": "generated"}
% BEGIN-VERIFY
% from sympy import symbols, exp, ln, Rational, solve, Eq, simplify
% t = symbols('t', positive=True)
% P = 200*exp(Rational(5,100)*t)
% sol = solve(Eq(P, 500), t)
% assert simplify(sol[0] - 20*ln(Rational(5,2))) == 0
% assert round(float(sol[0]), 2) == 18.33
% END-VERIFY
\begin{exercice}{TSPE-LOG-EX-003}{2}{15}
Une population de bacteries est modelisee par $P(t) = 200\,\mathrm{e}^{0{,}05t}$, ou $t$ est le temps en heures et $P(t)$ le nombre de bacteries (en milliers).
\begin{enumerate}
  \item Calculer $P(0)$. Interpreter.
  \item Determiner, en resolvant une equation, au bout de combien de temps la population atteint $500$ (milliers). Donner la valeur exacte puis une valeur approchee a $0{,}01$ heure pres.
\end{enumerate}
\end{exercice}
